\documentclass[11pt,a4paper]{article}
\usepackage[utf8]{inputenc}
\usepackage[T1]{fontenc}
\usepackage[swedish]{babel}
\usepackage{geometry}
\geometry{margin=2.5cm}
\usepackage{parskip}
% --- Preamble (Helvetica / sans) ---

\usepackage[scaled=0.92]{helvet}
\renewcommand{\familydefault}{\sfdefault}

\usepackage{microtype}
\usepackage{parskip}
\usepackage{tabularx}
\begin{document}

\begin{flushright}
Rickard Åberg \\
Malmö \\
\today
\end{flushright}

\vspace{1cm}

\section*{Personligt brev}

\textbf{Ansökan – Systemutvecklare} \\
\textbf{Referensnummer: 2026/15}

\vspace{0.5em}
Hej!

Mitt programmeringsintresse började med Basic på en Commodore 64 hemma framför tv:n och därefter assembler på Amiga500. Nyfikenheten på hur system fungerar har följt mig genom en civilingenjörsexamen i datateknik och över 15 års arbete med backendutveckling och systemintegration i större organisationer.

Redan som Problem Manager på Länstyrelsens IT avdelning kände jag att statliga myndigheter, det är där jag trivs bäst. 
Jag arbetar idag på Skogsstyrelsen ute i Skogarna vid \textit{Hemmestorps Mölla mellan Dalby och Blentarp i Skåne}, och har fått en första inblick i verksamheten. Jag blev nyfiken på er utvecklingsavdelning och därför söker jag härmed den här tjänsten som systemutvecklare.

Förutom min efarenhet av agil backend-utveckling har så är jag intresserad av AI och maskininlärning.  
Den solida grunden i Linjär Algebra, Flerdimensionell analys och statistiska proceser lades vid mina studier i Linköping och i år har jag kommit in på flertalet avancerade kurser samtidigt som jag arbetar ute i Skogen rent fysiskt.  Jag har redan påbörjat den första kursen i AI för naturlig språkbehandling vid Linköpings universitet och i nästa period läser jag en kurs i djup maskinlärning vid Växjö universitet (samt står som reserver på en kurs i Datorseende vid Halmstad samt Stora språkmodeller (LLM) vid Umeå).
\vspace{0.4em}
\noindent\textbf{Profil i korthet (nyckelord)}
\vspace{0.2em}

\noindent\begin{tabularx}{\linewidth}{@{}lX@{}}
\textbf{Backend} & Systemintegration, REST-tjänster, systemlandskap \\
\textbf{Arkitektur} & MicroServicedesign med fokus på förvaltningsbarhet \\
\textbf{Data} & SQL och databashantering i enterprise-miljöer \\
\textbf{Leverans} & Stor erfarenhet av DevOps, processer och stsabila flöden \\
\end{tabularx}

\vspace{0.4em}
Jag drivs av kvalitet och långsiktighet snarare än snabba genvägar, och jag trivs i team där samarbete, tydlig kommunikation och gemensamt ansvar är en självklar del av arbetssättet.

Som person är jag lugn, analytisk och idérik. Jag är social i samarbeten, men uppskattar också fokus och fördjupning. På fritiden dansar jag pardans, tränar, utövar yoga och tar kallbad – sådant som ger mig energi och balans. Hemma har jag min stora katt Diesel, som påminner mig om vikten av närvaro och att komma ut på en promenad för frisk luft.



Jag är bosatt i Malmö men både flyttbar, och van att arbeta på distans och ser fram emot att höra mer om denna möjlighet på en intervju!


Med vänlig hälsning
\textbf{Rickard Åberg}



\end{letter}
\end{document}
